\documentclass[10pt,a4paper,onecolumn]{article}

\usepackage{amsmath,amssymb}
\usepackage{graphicx}
\usepackage[margin=2cm]{geometry}
\usepackage{hyperref}
\usepackage{cite}
\usepackage{enumitem}
\usepackage{booktabs}

\hypersetup{colorlinks=true,linkcolor=black,urlcolor=blue,citecolor=blue}
\setlength{\parskip}{0.3ex}

\title{\textbf{Multi-Temporal Deforestation Segmentation \\
With Sentinel-2 and Global Forest Watch Labels} \\
\small Proof-of-Concept for Automated 10\,m Monitoring}

\author{Research Team --- Deforest.id}
\date{}

\begin{document}
\maketitle
\thispagestyle{empty}

\paragraph{Abstract.}
We present a proof-of-concept pipeline for automated deforestation
segmentation at 10\,m resolution using multi-temporal Sentinel-2
imagery and Global Forest Watch (GFW) lossyear labels.  Three
architectures are compared: a single-date U-Net (post-event only, 5
channels), a stacked multi-temporal U-Net (pre+post, 10 channels), and
a Siamese FC-Siam-diff network with explicit change features.  All
models are trained on paired 2020/2023 Sentinel-2 composites for
Indonesian protected forests.  The Siamese variant achieves the best
global test IoU of 0.137 (F1~0.240), outperforming both the stacked
multi-temporal (IoU~0.119) and single-date baseline (IoU~0.122).
Extreme class imbalance (0.38\,\% positive pixels) and the resolution
mismatch between GFW labels (30\,m Landsat) and Sentinel-2 imagery
(10\,m) limit absolute performance.  Despite modest metrics, the
pipeline demonstrates end-to-end viability: from satellite image
export through training to alert visualization.  We discuss
architectural choices, hyperparameter search results, failure modes,
and the path toward operational accuracy.

\section{Introduction}

Indonesia has the third-largest tropical forest area globally
($\sim$94~Mha), yet loses approximately 300~kha of natural forest
annually \cite{hansen2013high}.  Protected forests (\textit{hutan
lindung}) are especially vulnerable to small-scale deforestation that
evades coarse-resolution monitoring systems.  Operational products
such as Global Forest Watch (GFW) provide annual tree-cover loss at
30~m resolution, but this is insufficient for detecting clearings
below 0.9~ha and cannot support near-real-time interventions.

Deep learning on 10~m Sentinel-2 imagery offers a path toward more
timely and higher-resolution deforestation alerts.  Several recent
works apply U-Net variants to forest-change detection
\cite{wagner2019change, zapata2022deforestation}, typically using
single-date post-event imagery.  Multi-temporal approaches that fuse
pre- and post-event information should in principle outperform
single-date models by explicitly representing change, but practical
gains depend on the fusion architecture.

In this proof-of-concept study, we investigate three research questions:
\begin{enumerate}[nosep]
    \item Can GFW lossyear labels, despite their coarser 30\,m
          resolution, serve as viable training targets for 10\,m
          Sentinel-2 segmentation?
    \item Does explicit multi-temporal fusion (pre+post) improve
          over a single-date baseline for deforestation detection?
    \item What is the practical impact of fusion architecture ---
          naive channel stacking versus Siamese change detection?
\end{enumerate}

We build a complete pipeline covering data export, chip pairing,
oversampling, training with hyperparameter search, inference with
test-time augmentation, and metric aggregation.  All code and models
are publicly available at \url{https://github.com/anomalyco/deforest-id}.

\section{Related Work}

\paragraph{Deforestation monitoring.}
GFW \cite{hansen2013high} provides global annual forest-change data
at 30~m from Landsat.  Indonesia's SIMONTANA system uses a 6.25~ha
minimum mapping unit for REDD+ reporting.  Both are too coarse for
sub-hectare clearing detection.

\paragraph{Deep learning for change detection.}
U-Net \cite{ronneberger2015unet} is the dominant architecture for
biomedical and remote-sensing segmentation.  For change detection,
Daudt~et~al. \cite{daudt2018fcsiam} proposed FC-Siam-diff, a Siamese
encoder with shared weights and explicit feature differencing in the
decoder, which we adopt as our multi-temporal variant.

\paragraph{Class imbalance in deforestation.}
Deforestation pixels typically constitute less than 1\% of a scene.
Focal loss \cite{lin2017focal}, Dice loss \cite{milletari2016vnet},
and oversampling are standard remedies.  We find Dice loss with
10$\times$ positive-chip oversampling to be most effective.

\section{Methodology}

\subsection{Dataset}
We use Sentinel-2 L1C composite mosaics for two time periods:
\begin{itemize}[nosep]
    \item \textbf{Pre-event baseline:} 2020 cloud-free composite
          (5 bands: R, G, B, NIR, NDVI).
    \item \textbf{Post-event:} 2023 cloud-free composite
          (same 5 bands).
\end{itemize}
Both composites are tiled into $64\times64$ pixel chips at 10~m
resolution (0.64~ha per chip).  Pre- and post-event chips are paired
by spatial index.  Ground-truth masks are derived from the GFW
lossyear layer: pixels with lossyear~$\leq$~2023 are deforestation
(class~1), others are background (class~0).  Cloud-contaminated
pixels are masked via value~255 and ignored during loss computation.

Dataset statistics:
\begin{itemize}[nosep]
    \item Total paired chips: 40,484 (23,135 train, 5,822 val, 11,527 test).
    \item Global positive-pixel ratio: 0.38\,\%.
    \item Chips with zero deforestation: 74\,\%.
    \item Data sources: 13 Sentinel-2 scenes covering protected
          forests in Riau and Central Kalimantan.
\end{itemize}

\subsection{Architectures}

We evaluate three architectures:

\paragraph{Single-date baseline (SimpleUNet).}
A standard U-Net with 5 input channels (post-event R, G, B, NIR, NDVI
only).  Encoder uses ResNet-34-equivalent depth with 64 base filters,
21.3M parameters.  Trained with Dice loss for 50 epochs.  Best
validation IoU: 0.071.

\paragraph{Stacked multi-temporal (ResNet34 U-Net).}
Same U-Net architecture with 10 input channels (5 pre + 5 post
concatenated).  The first convolution layer is adapted from
ImageNet-pretrained ResNet34 weights.  21.3M parameters.  Best
validation IoU: 0.044 (epoch~5).

\paragraph{Siamese FC-Siam-diff.}
A shared ResNet34 encoder processes pre- and post-event chips
independently with tied weights.  The decoder fuses
$[\mathbf{f}_{\text{post}}, \mathbf{f}_{\text{pre}},
\mathbf{f}_{\text{post}} - \mathbf{f}_{\text{pre}}]$ at each
skip-connection level, explicitly modeling change.  25.9M parameters.
Best validation IoU: 0.054 (epoch~6).

\subsection{Training Procedure}

All multi-temporal models share:
\begin{itemize}[nosep]
    \item Optimizer: AdamW (lr=1e-4, weight decay=5e-4).
    \item Scheduler: Cosine annealing (T=30 epochs).
    \item Loss: Symmetric Dice loss (smooth=1.0).
    \item Batch size: 64 with gradient accumulation (effective 128).
    \item Oversampling: 10$\times$ weight for positive chips
          (contains deforestation).
    \item Augmentation: Random horizontal/vertical flips only.
    \item Precision: AMP (automatic mixed precision) with tf32.
    \item Gradient clipping: 1.0.
\end{itemize}

Initial experiments with ResNet18 (11M params), Tversky loss,
heavy augmentation (rotation, brightness jitter, Gaussian noise),
and larger weight decay all underperformed the above configuration.
Training runs 30 epochs ($\sim$15~minutes on RTX~3060 with RAM
caching) and is monitored via Rich TUI progress bars.

\section{Experiments and Results}

\subsection{Hyperparameter Search}

Nine configurations were systematically tested on the validation set:

\begin{table}[ht]
\centering
\caption{Hyperparameter search summary.  Config~9 is the best.}
\label{tab:hp_search}
\small
\begin{tabular}{lccl}
\toprule
Config & Encoder & Loss & Key Change \\
\midrule
1 & ResNet34 & Dice (symmetric)  & Baseline \\
2 & ResNet34 & Tversky ($\alpha$=0.7) & Recall boost \\
3 & ResNet34 & Focal ($\gamma$=2, $\alpha$=0.996) & Class balance \\
4 & ResNet34 & Lov\'asz-Softmax & Direct IoU opt. \\
5 & ResNet18 & Dice & Fewer params \\
6 & ResNet34 & Dice + accum=1 & No accumulation \\
7 & ResNet34 & Dice + accum=2 & Effective batch 128 \\
8 & ResNet34 & Dice + heavy aug & Rotation + noise \\
9 & ResNet34 & Dice + accum=2 + oversample 10$\times$ & \textbf{Best} \\
\bottomrule
\end{tabular}
\end{table}

Key findings:
\begin{itemize}[nosep]
    \item Dice loss consistently outperforms Tversky, Focal, and
          Lov\'asz for this task.
    \item Gradient accumulation (accum=2) was critical: without it,
          the model collapses to all-background prediction.
    \item 10$\times$ oversampling ensures $\sim$50\% of each batch
          contains positive chips.
    \item Heavy augmentation (rotation, brightness, noise) degrades
          performance --- the model needs consistent spectral patterns.
    \item ResNet34 (21M params) outperforms ResNet18 (11M).
\end{itemize}

\subsection{Test-Set Results}

The three best architectures are evaluated on 11,390 hold-out test
chips with test-time augmentation (horizontal + vertical flip
ensemble).

\begin{table}[ht]
\centering
\caption{Test-set comparison: global metrics (summed TP/FP/FN across all chips).}
\label{tab:results}
\begin{tabular}{lccc}
\toprule
Metric & Single-date (post) & Multi-temporal (stacked) & Siamese (ours) \\
\midrule
Global IoU   & 0.122 & 0.119 & \textbf{0.137} \\
Precision    & \textbf{0.198} & 0.154 & 0.196 \\
Recall       & 0.239 & \textbf{0.344} & 0.311 \\
F1           & 0.217 & 0.212 & \textbf{0.240} \\
Accuracy     & \textbf{0.984} & 0.977 & 0.982 \\
\bottomrule
\end{tabular}
\end{table}

The Siamese model achieves the best IoU and F1, combining competitive
precision (0.196) with strong recall (0.311).  The stacked
multi-temporal model has the highest recall but lowest precision ---
it detects more deforestation but at the cost of false positives.
The single-date baseline has the best precision but misses many
events.

Per-sample mean IoU is substantially lower (0.012--0.015) because
74\,\% of test chips contain zero deforestation (IoU=0).  Global
(summed) metrics better reflect true pixel-level performance.

\subsection{Overfitting and Generalization}

All multi-temporal configurations show peak validation IoU at
epochs~4--6 followed by gradual decline (Figure~\ref{fig:curves}).
Training loss continues to decrease while validation loss plateaus
or increases, indicating overfitting.  This pattern is consistent
across all 9 configurations, suggesting the 64$\times$64 chip size
provides limited training signal for 21M+ parameter models.

\begin{figure}[ht]
\centering
\includegraphics[width=0.9\textwidth]{../reports/training_curves.png}
\caption{Training curves for the best Siamese configuration:
loss, IoU, precision/recall/F1, and learning rate over 30 epochs.
Overfitting begins after epoch~6.}
\label{fig:curves}
\end{figure}

\subsection{Error Analysis}

Qualitative examination of 20 random test samples reveals:
\begin{enumerate}[nosep]
    \item \textbf{False positives} cluster around spectral changes
          from non-forest transitions (agriculture, bare soil, water
          level changes).
    \item \textbf{Boundary over-smoothing}: the model misses small
          clearings ($<$0.5~ha) and produces overly smooth edges.
    \item \textbf{Cloud-shadow residuals} in the composites produce
          consistent false positives despite cloud masking.
    \item \textbf{Label noise}: many GFW deforestation polygons
          align poorly with visible spectral change in 10~m imagery,
          indicating the 30~m Landsat origin introduces spatial
          uncertainty.
\end{enumerate}

\section{Proof-of-Concept Implementation}

The complete pipeline runs on a single consumer GPU.  Key design
decisions:

\paragraph{Data pipeline.}
Sentinel-2 composites are exported as GeoTIFF via Google Earth Engine,
then tiled into 64$\times$64 NPZ chips (5 spectral bands + cloud mask).
Pre/post pairing is done by spatial index matching.  All chips are
cached to RAM at training start for I/O-free throughput ($\sim$27
sec/epoch on 23k chips).

\paragraph{Training.}
Implemented in PyTorch 2.7 with AMP and tf32 matrix precision.
Training survives SSH disconnect via tmux.  The Rich TUI library
provides live progress bars and metric tables.  CSV logging enables
post-hoc analysis.

\paragraph{Inference.}
Test-time augmentation flips each chip horizontally and vertically,
averaging the three predictions.  Per-chip visualizations
(6-panel: pre-RGB, post-RGB, NDVI pre/post, GT mask, prediction,
overlay) are saved for qualitative inspection.

\paragraph{Hardware.}
All experiments run on NVIDIA RTX~3060 (12~GB VRAM) within WSL2
(Windows Subsystem for Linux).  Training completes in under 30
minutes for 30 epochs.  Full test-set inference (11,390 chips with
TTA) requires approximately 12~minutes.

\section{Limitations and Challenges}

Despite demonstrating a working pipeline, several limitations prevent
operational deployment:

\paragraph{GFW label noise.}
The most significant challenge.  GFW lossyear data originates from
Landsat (30~m), while our composites are Sentinel-2 (10~m).  This
3$\times$ resolution mismatch means:
\begin{itemize}[nosep]
    \item A single GFW pixel covers nine Sentinel-2 pixels, creating
          ambiguous supervision at the 10~m scale.
    \item Deforestation timing mismatch: GFW reports annual loss,
          while our composites capture specific pre/post dates,
          introducing temporal label noise.
    \item Small clearings visible in 10~m imagery may be absent from
          the 30~m GFW label.
\end{itemize}

\paragraph{Extreme class imbalance.}
With only 0.38\,\% positive pixels and 74\,\% of chips containing
zero deforestation, the model has limited training signal.
Oversampling helps but cannot create new information --- the model
sees the same few positive chips repeatedly.

\paragraph{Chip size vs.\ context.}
The 64$\times$64 chip (0.64~ha) restricts spatial context.  Larger
chips (128$\times$128 or 256$\times$256) would provide more
boundary information but reduce the effective number of training
samples.

\paragraph{Compositing artifacts.}
Cloud-free composites average multiple acquisitions, which can blur
deforestation events that occurred mid-period.  Temporal
misalignment between pre/post composites also introduces false
change signals.

\paragraph{Distribution shift.}
The test set has higher deforestation density (0.9\,\%) than the
training set (0.38\,\%), indicating spatial heterogeneity that
challenges generalization.

\section{Conclusion and Future Work}

We presented a proof-of-concept pipeline for multi-temporal
deforestation segmentation at 10~m resolution.  The Siamese
FC-Siam-diff architecture achieves the best result (IoU~0.137,
F1~0.240), confirming that explicit change representations outperform
naive channel stacking.  Despite modest absolute metrics, the
pipeline demonstrates end-to-end viability: from GEE data export to
trained model with inference visualization, all on a single consumer
GPU.

Key lessons for practitioners:
\begin{itemize}[nosep]
    \item GFW labels are usable but noisy at 10~m --- better
          training data would substantially improve results.
    \item Dice loss + aggressive oversampling + gradient
          accumulation are necessary for extreme imbalance.
    \item Multi-temporal fusion architecture matters more than
          whether fusion is used at all.
\end{itemize}

Future work includes:
\begin{enumerate}[nosep]
    \item Higher-quality labels via manual annotation or
          high-resolution imagery (Planet, 3~m).
    \item Larger chip sizes (256$\times$256) with attention
          mechanisms.
    \item Sentinel-1 SAR integration for cloud-free change detection.
    \item Post-processing with connected-component filtering and
          CRF refinement.
\end{enumerate}

\small
\bibliographystyle{ieeetr}
\bibliography{icp-refs}

\end{document}
